\documentclass[11pt,a4paper]{article}
\usepackage[margin=25mm]{geometry}
\usepackage{amsmath,amssymb,amsthm,hyperref}
\setlength{\emergencystretch}{2em}
\newtheorem{proposition}{Proposition}
\newcommand{\gen}[1]{\langle#1\rangle}
\title{A transitive tail beyond a singleton endpoint}
\author{Human Compression project: bounded L044 proof component}
\date{}
\begin{document}\maketitle
Let $G=\gen{\rho_0,\ldots,\rho_{r-1}}\leq A_n$ be a faithful
transitive string C-group, $n\geq12$, with every maximal
omission intransitive. Suppose $G_0$ has a singleton orbit
$\{a\}$, its other orbit contains $b$, and $ab$ is the unique
eligible $0$-edge. Assume $K=G_{>1}$ is transitive on
$X=\Omega\setminus\{a,b\}$ and fixes $a,b$.

\begin{proposition}
Under these hypotheses, $r\leq\lfloor(n-1)/2\rfloor$.
\end{proposition}
\begin{proof}
The crossing involution has the form
$\rho_0=(a\ b)\lambda$, with $\lambda$ an odd involution on
$X$. It centralizes the transitive action of $K$ on $X$.
Therefore $\lambda$ has no fixed point: a fixed point of a
centralizing permutation would propagate to the entire orbit.
Write $|X|=2m$. Then $\lambda$ consists of $m$ transpositions,
so $m$ is odd. As $n=2m+2\geq12$, one has $m\geq5$.

The pairs of $\lambda$ are a $K$-invariant block system.
Every distinguished $K$ generator is essential in the block
action. Indeed choose a minimal distinguished subset producing
the whole block image. It has at most two point orbits on $X$;
adjoining $\rho_0$, whose restriction is $\lambda$, joins
them. If a $K$ generator were omitted by this subset, the
corresponding global omission would be transitive on $X$.
That omitted generator is supported on $X$, so adjoining it
cannot merge omission orbits. Since the result is transitive
$G$, the omission would already be transitive, a contradiction.
This is the full actual-omission argument of the accepted
commuting-tail lemma, specialized to the present pair system.

Whiston's bound now gives $r-2\leq m-1$, hence $r\leq m+1$.
As $\lfloor(n-1)/2\rfloor=m$, a failure of the desired bound
would force $r=m+1$ and make the $m-1$ independent block images
generate the full $S_m$, by Whiston's equality clause.

Put $\sigma=\rho_1$. It fixes $a$, because it belongs to $G_0$,
and it must move $b$, because $K$ fixes $b$ whereas
$G_0=\gen{\sigma,K}$ is transitive on $\{b\}\cup X$.
Thus for some $x\in X$ one has
\[
                       \sigma=(b\ x)v,
\]
where $v$ is an odd involution supported on $X\setminus\{x\}$.
Let $x'=x^\lambda$ and $H=G_{>2}$. All its generators fix $b$
and commute with $\sigma$, so fix $x$. They also commute with
$\lambda$, so fix $x'$. Their $m-2$ block images are independent
and fix the pair $\{x,x'\}$; Whiston makes their image the full
$S_{m-1}$ on the other pairs.

The only $H$-fixed points in $X$ are $x,x'$, since its block
image moves every other pair. As $v$ centralizes $H$ and fixes
$x$, it fixes $x'$ as well. On the remaining $m-1$ pairs,
the accepted paired-row centralizer lemma applies: their
number is even and at least four, and the induced group is
$S_{m-1}$. Hence $v$ is even on the remaining points.
This contradicts the oddness of $v$, excludes $r=m+1$ and
proves the proposition.
\end{proof}

\paragraph{Source and scope.}
This is a complete replacement proof of the transitive-tail
singleton-endpoint case of Cameron--Fernandes--Leemans--Mixer,
\emph{Highest rank of a polytope for $A_n$},
arXiv:1605.09173v1, Proposition~5.29, physical page~38, in the
main $n\geq12$ range. It reuses the accepted commuting-tail
and paired-row lemmas and retained Whiston bounds. The tight
case is excluded without importing the unproved full
equality-graph assertion of source Proposition~5.21.
No finite computation, projected IP or new formal theorem
is required. The intransitive tail, the distinct $D_{10}$
residual, later table exhaustions and the complete L044
source remain separate. No original or question credit is
asserted.
\end{document}
